\documentclass[a4paper,12pt]{article}

\input{../../preamble_problems}

\begin{document}

\begin{center}
\Large{Tutorial Exercises - Week 1.1} \\
\vspace{3mm}
\large{4001CMD - Mathematical Skills for Computing Professionals} \\
\vspace{5mm}
\end{center}

The following exercises should be attempted before your next tutorial session this week. They focus on material covered in Week 1, namely set theory and symbolic logic. 

\section*{Set Theory}

\begin{problem}
	For the following, show rigorously that $A = B$. 
	\begin{enumerate}[label=(\alph*)]
		\item $A = \{1,2,3\}$, $B = \{n : n \in \N, n^2 < 10\}$. 
		\item $A = \{x \in \R: x^2 - 4x + 4 = 1\}$, $B = \{1,3\}$
	\end{enumerate}
\end{problem}

\begin{solution}
	(a) $1^2 < 10$, $2^2 = 4 < 10$, and $3^2 = 9 < 10$, so $A \subseteq B$. Conversely, if $n \in B$, then $n \in \N$, and $n^2 < 10$. There are only three numbers with this property: 1, 2, and 3, so $n \in A$ in this case. So $B \subseteq A$. \qed
	
	(b) If $x \in A$, then $x^2 - 4x + 4 = 1$, so $x^2 - 4x + 3 = 0$. This factors to $(x - 3)(x-1) = 0$, so either $x=1$ or $x=3$, so $x \in B$. Conversely, if $x \in B$, then $x = 1$ or $x=3$; in either case, $x^2 - 4x + 4 = 1$ is satisfied, so $x \in A$, so $B \subseteq A$. \qed
\end{solution}

\begin{problem}
	For the following, show rigorously that $A \neq B$. 
	\begin{enumerate}[label=(\alph*)]
		\item $A = \{1,2\}$, $B=\{x \in \R : x^3 - 2x^2 - x + 2 = 0\}$
		\item $A = \{1,3,5\}$, $B = \{n : n \in \N, n^2 - 1 \leq n\}$
	\end{enumerate}
\end{problem}

\begin{solution}
	(a) The polynomial factors as $x^3 - 2x^2 - x + 2 = x^2(x-2) - (x-2) = (x^2 - 1)(x-2) = (x-1)(x+1)(x-2)$. So while $A \subseteq B$, the number $x = -1$ lives in $B$, but $-1 \not\in A$, so $B \neq A$. \qed
	
	(b) For $n = 3$, we have $n^2 - 1 = 8 > 3$, so $3 \not\in B$, so $A \neq B$. \qed 
\end{solution}

\begin{problem}
	For the following, determine whether $A = B$. 
	\begin{enumerate}[label=(\alph*)]
		\item $A = \{x : x^2 + x = 2\}$, $B = \{1,-1\}$
		\item $A = \{x \in \R : 0 < x \leq 2\}$, $B = \{1,2\}$
	\end{enumerate}
\end{problem}

\begin{solution}
	(a) For $x=-1$, we have $x^2 + x = 1 - 1 = 0 \neq 2$, so $-1 \not\in A$, so $B \neq A$. 
	
	(b) The number $x = 1/2$ lies between 0 and 2, but $1/2 \not\in B$, so $A \neq B$. 
\end{solution}

\begin{problem}
	For the following, determine whether that $A \subseteq B$. 
	\begin{enumerate}[label=(\alph*)]
		\item $A = \{1,2\}$, $B = \{x : x^3 - 6x^2 + 11x = 6\}$
		\item $A = \{2n : n \in \N\}$, $B = \N$
		\item $A = \{x \in \R : x^3 - 2x - x + 2 = 0\}$, $B =\{1,2\}$
		\item $A = \{1,2,3,4\}$, $C = \{5,6,7,8\}$, $B = \{n \in \N : n+m = 8 \textrm{ for some } m \in C\}$
	\end{enumerate} 
\end{problem}

\begin{solution}
	(a) Plugging in $x=1$ and $x=2$ into the polynomial in the definition of $B$, we get $1 - 6 + 11 = -5+11 = 6$, and $8 - 24 + 22 = 8-2 = 6$, so $1 \in B$ and $2 \in B$, so $A \subseteq B$. 
	
	(b) Since $2n$ is a positive integer for every positive integer $n$, we have $2n \in \N$ for all $n \in \N$, so $A \subseteq B$. 
	
	(c) As we saw in Problem 1(a), $A \not\subseteq B$, since $-1 \in A$ but $-1 \not\in \{1,2\}$. 
	
	(d) Note that for $4 \in A$, nothing in $C$ sums with $4$ to give $8$, so $4 \in A$ but $4 \not\in B$, so $A \not\subseteq B$. 
\end{solution}

\begin{problem}
	The \emph{symmetric difference} fo two sets $A$ and $B$ is the set defined by \[
	A \triangle B = (A \cup B) \setminus (A \cap B)
	\]
	(recall for sets $X$ and $Y$ with $Y \subseteq X$, $X \setminus Y$, read ``$X$ minus $Y$'', is the set of elements of $X$ that are not elements of $Y$). 
	\begin{enumerate}[label=(\alph*)]
		\item If $A = \{1,2,3\}$ and $B = \{2,3,4,5\}$, find $A \triangle B$. 
		\item Describe in words the symmetric difference of two sets $A$ and $B$ (not necessarily those sets in part (a)). 
	\end{enumerate} 
\end{problem}

\begin{solution}
	(a) $A \triangle B = \{1,2,3,4,5\} \setminus \{2,3\} = \{1,4,5\}$ -- elements in the Union, which are not in the intersection

	(b) $A \triangle B$ are those sets that are in either $A$ or in $B$, but not both. (This is the ``exclusive or'' of set theory A "XOR" B.)
\end{solution}

\section*{Symbolic Logic}

\begin{problem}
	Given that $p$ is false, $q$ is true, $r$ is false, and $s$ is true, determine whether the following propositions are true or false. 
	\begin{enumerate}[label=(\alph*)]
		\item $\neg p \vee \neg(q \wedge r)$
		\item $\neg(p \vee q) \wedge (\neg p \vee r)$
		\item $(p \vee \neg r) \wedge \neg ((q \vee r) \vee \neg (r \vee p))$
		\item $(p \then q) \wedge (q \then r)$
		\item $(p \then q) \then r$
		\item $p \then (q \then r)$
		\item $(s \then( p \wedge \neg r)) \wedge ((p \then (r \vee q)) \wedge s)$
		\item $((p \wedge \neg q) \then (q \wedge r)) \then (s \vee \neg q)$
	\end{enumerate} 
\end{problem}

\begin{solution}
	(a) TRUE: because $\neg p$ is true. 
	
	(b) FALSE: $p \vee q$ is true, so $\neg(p \vee q)$ is false, so the proposition is false. 
	
	(c) FALSE: $q \vee r$ is true since $q$ is true, so $(q \vee r) \vee \neg(r \vee p)$ is true. This means $\neg((q \vee r) \vee \neg(r \vee p))$ is false. Since the main expression is an AND statement, the statement is false.
	
	(d) FALSE: Since $q$ is true but $r$ is false, $q \then r$ is false, so the proposition is false. 
	
	(e) FALSE: Since $p$ is false, $p \then q$ is true. But $r$ is false. So the proposition $(p \then q) \then r$ is false. 
	
	(f) TRUE: Since $p$ is false, the proposition is true regardless of whether $q \then r$ is true (it's not). 
	
	(g) FALSE: Since $p \wedge \neg r$ = $F \wedge T$ is false, and  $s$ is true, so $s \then (p \wedge \neg r)$ is $T \then F=F$. Also,  $p \then (r \vee q)$ is $F \then (F \vee T)=$ true, and so is $s$, hence right part is $T\wedge T = T$. So the whole proposition is then $F \wedge T=F$. 
	
	(h) TRUE: Since $p$ is false, $p \wedge \neg q$ is false, so the proposition $(p \wedge \neg q) \then (q \wedge r)$ is true. Meanwhile, since $s$ is true, $s \vee \neg q$ is true, so the proposition is true. 
\end{solution}

\begin{problem}
	For each of the following, formulate the symbolic expression in words, assuming the propositions stand for the following statements: 
	\begin{align*}
	p &: \textrm{ You play football.} \\
	q &: \textrm{ You miss the midterm exam.} \\
	r &: \textrm{ You pass the module.}
	\end{align*}
	\begin{enumerate}[label=(\alph*)]
		\item $p \vee q \vee r$
		\item $\neg(p \vee q) \vee r$
		\item $(p \wedge q) \vee (\neg q \wedge r)$
	\end{enumerate} 
\end{problem}

\begin{solution}
	(a) You play football, or you miss the midterm exam, or you pass the module. 
	
	(b) You neither play football nor miss the midterm exam, or you pass the module. 
	
	(c) You play football and miss the midterm exam, or you be there for the midterm exam and you pass the module. 
\end{solution}

\begin{problem}
	For each of the following, represent the proposition symbolically by letting
	\begin{align*}
	p &: \textrm{ You run 10 laps daily.} \\
	q &: \textrm{ You are healthy.} \\
	r &: \textrm{ You take multi-vitamins.}
	\end{align*}
	\begin{enumerate}[label= (\alph*)]
		\item You run 10 laps daily, but you are not healthy. 
		\item You do not run 10 laps daily, you do not take multi-vitamins, and you are not healthy. 
		\item Either you are healthy or you do not run 10 laps daily, and you do not take multi-vitamins. 
	\end{enumerate}
\end{problem}

\begin{solution}
	\vspace{2mm}
	
	(a) $p \wedge \neg q$. 
	
	(b) $(\neg p )\wedge (\neg q) \wedge (\neg r)$
	
	(c) $(q \vee \neg p) \wedge \neg r$
\end{solution}

%\begin{problem}
%	Formulate the following arguments symbolically, and determine whether each is valid, using: 
%	\[
%	p : \textrm{ I study hard.} \quad q : \textrm{ I get A's.} \quad r : \textrm{ I get rich.}
%	\]
%	\begin{enumerate}[label=(\alph*)]
%	\item
%	\begin{tabular}{l}
%		If I study hard, then I get A's. \\
%		I study hard. \\
%		\hline
%		$\therefore$ I get A's.
%	\end{tabular} 
%	\item 
%	\begin{tabular}{l}
%		If I study hard, then I get A's. \\
%		If I don't get rich, then I don't get A's. \\
%		\hline
%		$\therefore$ I get rich. 
%	\end{tabular} 
%	\item \begin{tabular}{l}
%		I study hard if and only if I get rich. \\
%		I get rich. \\
%		\hline
%		$\therefore$ I study hard. 
%	\end{tabular} 
%	\item \begin{tabular}{l}
%		If I study hard or I get rich, then I get A's. \\
%		 I get A's. \\
%		 \hline
%		 $\therefore$ If I don't study hard, then I get rich. 
%	\end{tabular} 
%	\end{enumerate}
%\end{problem}
%
%\begin{solution}
%	(a) \begin{center} 
%		\begin{tabular}{l}
%		$p \then q$ \\
%		$p$ \\
%		\hline
%		$\therefore q$
%	\end{tabular}
%	\end{center} 
%	This is valid by \emph{modus ponens}. 
%	
%	(b)
%	\begin{center}
%		\begin{tabular}{l}
%			$p \then q$ \\
%			$\neg r \then \neg q$ \\
%			\hline
%			$\therefore r$
%		\end{tabular} 
%	\end{center}
%	This is not valid. The second statement is equivalent to $q \then r$; by hypothetical syllogism and the first statement, this implies $p \then r$. But since neither $p$ nor $q$ are assumed, we cannot conclude $r$. 
%	
%	(c)
%	\begin{center}
%		\begin{tabular}{l}
%			$p \Leftrightarrow r$ \\
%			$r$ \\
%			\hline
%			$\therefore p$
%		\end{tabular} 
%	\end{center}
%	This is valid: 
%	\begin{center}
%		\begin{tabular}{llr}
%			1. & $p \Leftrightarrow r$ & Assumption 1 \\
%			2. & $(p \then r) \wedge (r \then p)$ & Definition of $\Leftrightarrow$ \\
%			3. & $(r \then p) \wedge (p \then r)$ & Commutativity of $\wedge$ \\
%			4. & $r \then p$ & Simplification, line 3. \\
%			5. & $r$ & Assumption 2 \\
%			6. & $p$ & \textit{Modus ponens}, lines 4 and 5. 
%		\end{tabular} 
%	\end{center}
%	
%	(d) 
%	\begin{center}
%		\begin{tabular}{l}
%			$(p \vee r) \then q$ \\
%			$q$ \\
%			\hline
%			$\therefore (\neg p) \then r$
%		\end{tabular} 
%	\end{center}
%	This is not valid. Even if $(p \vee r) \then q$ is true, the truth value of $q$ tells us nothing about the truth value of $p \vee r$, and thus, tells us nothing about the truth value of $p$ or $r$ individually or in relation to each other. So we cannot conclude that $(\neg p) \then r$ is true. 
%\end{solution}
%
%\begin{problem}
%	For each of the following, write the given argument in words and determine whether each argument is valid. Let: \begin{align*}
%	p &: \textrm{ The for loop is faulty.} \\
%	q &: \textrm{ The while loop is faulty.} \\
%	r &: \textrm{ The hardware is unreliable.} \\
%	s &: \textrm{ The output is correct.} 
%	\end{align*}
%	\begin{enumerate}[label=(\alph*)]
%		\item \begin{tabular}{l}
%			$(r \vee s) \then p$ \\
%			$s \then q$ \\
%			$p \vee s$ \\
%			\hline
%			$\therefore r$
%		\end{tabular} 
%		\item \begin{tabular}{l}
%			$(p \then r) \then q$ \\
%			$(q \then s) \then p$ \\
%			$r \wedge s$ \\
%			$p \vee q$ \\
%			\hline
%			$\therefore p \wedge q$
%		\end{tabular} 
%		\item \begin{tabular}{l}
%			$p \then q$ \\
%			$q \then r$ \\
%			$\neg r$ \\
%			$s \then r$ \\
%			\hline 
%			$\therefore \neg s$
%		\end{tabular} 
%		\item \begin{tabular}{l}
%			$p \then (q \vee r)$\\
%			$q \then (p \vee s)$ \\
%			$p \vee \neg q$ \\
%			$\neg s$ \\
%			\hline
%			$\therefore p \vee r$
%		\end{tabular}
%	\end{enumerate} 
%\end{problem}
%
%\begin{solution}
%	(a) If either the hardware is unreliable or the output is correct, then the for loop is faulty. If the output is correct, then the while loop is faulty. Either the for loop is faulty, or the output is correct. So the hardware is unreliable. 
%	
%	This is not valid. We know nothing about $p$ on its own, so we cannot conclude whether $r \vee s$ is true or false, and we cannot conclude anything about $s$ either. 
%	
%	(b) If the hardware is unreliable when the for loop is faulty, then the while loop is faulty. If the output is correct when the while loop is faulty, then the for loop is faulty. The hardware is unreliable, but the output is correct. And either the for loop or the while loop is faulty. So actually both the for loop and the while loop are faulty. 
%	
%	This is valid: 
%	\begin{center}
%		\begin{tabular}{llr}
%			1. & $(p \then r) \then q$ & Assumption 1 \\
%			2. & $(q \then s) \then p$ & Assumption 2 \\
%			3. & $r \wedge s$ & Assumption 3 \\
%			4. & $p \vee q$ & Assumption 4 \\
%			5. & $r$ & Simplification, line 3 \\
%			6. & $s \wedge r$ & Commutativity, line 3 \\
%			7. & $s$ & Simplification, line 6 \\
%			8. & $p \then r$ & Definition of $\then$, line 5 \\
%			9. & $q \then s$ & Definition of $\then$, line 7 \\
%			10. & $q$ & \textit{Modus ponens}, lines 1 and 8 \\
%			11. & $p$ & \textit{Modus ponens}, lines 2 and 9 \\
%			12. & $p \wedge q$ & Conjunction, lines 11 and 10. 
%		\end{tabular} 
%	\end{center}
%	
%	(c) The while loop is faulty when the for loop si faulty. The hardware is unreliable when the while loop is faulty. But the hardware is reliable. If the output is correct, then the hardware is unreliable. So the output is incorrect. 
%	
%	This is valid: 
%	\begin{center}
%		\begin{tabular}{llr}
%			1. & $\neg r$ & Assumption 3 \\
%			2. & $s \then r$ & Assumption 4 \\
%			3. & $\neg s$ & \textit{Modus tolens}, lines 2 and 1. 
%		\end{tabular} 
%	\end{center}
%	
%	(d) If the for loop is faulty, then either the while loop is faulty or the hardware is unreliable. If the while loop is faulty, then either the for loop is faulty or the output is correct. Either the for loop is faulty, or the while loop is functioning. The output is not correct. Therefore either the for loop is faulty or the hardware is unreliable. 
%	
%	This is not valid. If $p$ is false, then we cannot conclude that $r$ is true, and if $r$ is false, we cannot conclude that $p$ is true (we could if $q$ were false, by assumption 1). 
%\end{solution}
%
%\begin{problem}
%	Write a symbolic logic expression equivalent to the following logic circuits. 
%	\begin{enumerate}[label=(\alph*)]
%		\item \ \\ \begin{center}
%			\includegraphics[scale=0.3]{circuit1}
%		\end{center}
%		\item \ \\ \begin{center}
%			\includegraphics[scale=0.3]{circuit2}
%		\end{center}
%		\item \ \\ \begin{center}
%			\includegraphics[scale=0.3]{circuit3}
%		\end{center}
%		\item (This one is hard!)
%		\begin{center}
%			\includegraphics[scale=0.3]{circuit5}
%		\end{center}
%	\end{enumerate}
%\end{problem}
%
%\begin{solution}
%	(a) $\neg(x_1 \wedge x_2)$
%	
%	(b) $(x_1 \wedge  x_2) \vee \neg x_3$
%	
%	(c) $\neg((\neg x_1) \vee x_2 )$
%	
%	(d) $(((\neg x_1) \vee x_2) \wedge ((\neg x_3) \vee x_4)) \wedge ((\neg x_2) \vee x_4)$
%\end{solution}
%
%\begin{problem}
%	Show that the following logic circuits are equivalent.  
%	\begin{enumerate}[label=(\alph*)]
%		\item \ \\ \begin{center}
%			\includegraphics[scale=0.3]{circuits1}
%		\end{center}
%		\item \ \\ \begin{center}
%			\includegraphics[scale=0.3]{circuits2}
%		\end{center}
%		\item \ \\ \begin{center}
%			\includegraphics[scale=0.3]{circuits3}
%		\end{center}
%		\item \ \\
%		\begin{center}
%			\includegraphics[scale=0.3]{circuits4}
%		\end{center}
%	\end{enumerate}
%\end{problem}
%
%\begin{solution}
%	(a) The logic circuits are $\neg(x_1 \wedge x_2)$ and $(\neg x_1) \vee (\neg x_2)$. These are equivalent by de Morgan's laws.
%	
%	(b) The logic circuits are expressible as $(\neg(x_1 \vee (\neg x_2))) \vee ((\neg x_1) \wedge x_3)$ and $(\neg x_1) \wedge (x_2 \vee x_3)$. By de Morgan's law, $\neg(x_1 \vee(\neg x_2))$ is equivalent to $(\neg x_1) \wedge x_2$. So the statement is expressible as $((\neg x_1) \wedge x_2) \vee ((\neg x_1) \wedge x_3)$. By distributivity, this is $(\neg x_1) \wedge (x_2 \vee x_3)$. 
%	\begin{center}
%		\begin{tabular}{llr}
%			1. & $(\neg(x_1 \vee (\neg x_2))) \vee ((\neg x_1) \wedge x_3)$ & Initial circuit \\
%			2. & $((\neg x_1) \wedge x_2) \vee ((\neg x_1) \wedge x_3)$ & de Morgan's law applied to $\neg(x_1 \vee (\neg x_2))$ \\
%			3. & $(\neg x_1) \wedge (x_2 \vee x_3) $ & Distributivity, line 2 
%		\end{tabular} 
%	\end{center} 
%	
%	(c) The logic circuits are expressible as $(x_1 \wedge x_2) \vee x_1$ and $(x_1 \vee x_2) \wedge x_1$. We have: 
%	\begin{center}
%		\begin{tabular}{llr}
%			1. & $(x_1 \wedge x_2) \vee x_1$ & Initial circuit \\
%			2. & $x_1 \vee (x_1 \wedge x_2)$ & Commutativity, line 1 \\
%			3. & $(x_1 \vee x_1) \wedge (x_1 \vee x_2)$ & Distributivity, line 2 \\
%			4. & $(x_1 \vee x_1) \Leftrightarrow x_1 $ & Tautology \\
%			5. & $x_1 \wedge (x_1 \vee x_2)$ & Substitution, lines 3 and 4
%		\end{tabular} 
%	\end{center}
%	
%	(d) The circuits are $(\neg x_1) \vee ((\neg x_2) \vee x_3)$ and $(\neg (x_1 \wedge x_2)) \vee x_3$. We have: 
%	\begin{center}
%		\begin{tabular}{llr}
%			1. & $(\neg x_1) \vee ((\neg x_2) \vee x_3)$ & Initial circuit \\
%			2. & $((\neg x_1) \vee (\neg x_2)) \vee x_3$ & Associativity, line 1 \\
%			3. & $(\neg (x_1 \wedge x_2)) \vee x_3$ & De Morgan's law, line 2
%		\end{tabular} 
%	\end{center}
%\end{solution}
%\begin{itemize}
%	\item Addition: 
%	\begin{center}
%		\begin{tabular}{r}
%			$P$ \\
%			\hline
%			$\therefore P \vee Q$
%		\end{tabular}
%	\end{center}
%	\item Conjunction: 
%	\begin{center}
%		\begin{tabular}{r}
%			$P$ \\
%			$Q$ \\
%			\hline
%			$\therefore P \wedge Q$
%		\end{tabular}
%	\end{center}
%	\item Simplification: 
%	\begin{center}
%		\begin{tabular}{r}
%			$P \wedge Q$ \\
%			\hline
%			$\therefore P$
%		\end{tabular}
%	\end{center}
%	\item Disjunctive syllogism: 
%	\begin{center}
%		\begin{tabular}{r}
%			$P \vee Q$ \\
%			$\neg P$ \\
%			\hline
%			$\therefore Q$
%		\end{tabular}
%	\end{center}
%	\item Hypothetical syllogism:
%	\begin{center}
%		\begin{tabular}{r}
%			$P \then Q$ \\
%			$Q \then R$ \\
%			\hline
%			$\therefore P \then R$
%		\end{tabular}
%	\end{center}
%	\item \textit{Modus ponens}:
%	\begin{center}
%		\begin{tabular}{r}
%			$P \then Q$ \\
%			$P$ \\
%			\hline
%			$\therefore Q$
%		\end{tabular}
%	\end{center}
%	\item \textit{Modus tollens}:
%	\begin{center}
%		\begin{tabular}{r}
%			$P \then Q$ \\
%			$\neg Q$ \\
%			\hline
%			$\therefore \neg P$
%		\end{tabular}
%	\end{center}
%	\item Commutativity: 
%	\[
%	P \vee Q \iff Q \vee P \qquad P \wedge Q \iff Q \wedge P
%	\]
%	\item Associativity: 
%	\[
%	P \vee (Q \vee R) \iff (P \vee Q) \vee R \qquad P \wedge (Q \wedge R) \iff (P \wedge Q) \wedge R
%	\]
%	\item de Morgan's Laws:
%	\[
%	\neg (P \vee Q) \iff (\neg P) \wedge (\neg Q) \qquad \neg(P \wedge Q) \iff (\neg P) \vee(\neg Q)
%	\]
%	\item Negation: 
%	\[
%	P \iff \neg\neg P \qquad P \vee \neg P \textrm{ is TRUE} \qquad P \wedge \neg P \textrm{ is FALSE}
%	\]
%	\item Tautology: 
%	\[
%	P \iff (P \vee P) \qquad P \iff (P \wedge P)
%	\]
%\end{itemize}
%
%Additionally, you may use Substitution: if $X \iff Y$ is true, and $X$ appears in a proposition, you may replace $X$ with $Y$. 


\end{document}